
\subsection*{3.2}
\phantomsection
\label{sga3:VIIB:3.2}

Suppose now that \(C\) is Artinian, with radical \(\mathfrak r\).
For every \(C\)-algebra \(U\) \textup{(associative and unital)}, the
ideal \(\mathfrak rU\) therefore consists of nilpotent elements. If
\(x\in\mathfrak rU\), put
\[
  \exp_U(x)
  =
  1+x+\frac{x^2}{2!}+\cdots.
\]
This gives a bijection from \(\mathfrak rU\) onto
\(1+\mathfrak rU\); its inverse sends \(1+y\in1+\mathfrak rU\) to
\[
  \log_U(1+y)
  =
  y-\frac{y^2}{2}+\frac{y^3}{3}-\cdots.
\]
Moreover, the map \(\exp_U\) is plainly functorial in
\(U\).

Still assuming \(C\) Artinian, suppose that \(U\) is endowed with a
bialgebra structure over \(C\) \textup{(cf.\ \SGARefLink{sga3:VIIA:2.2}{2.2})}. For every
primitive element \(x\in\mathfrak rU\)
\textup{(cf.\ VII\textsubscript{A}, \SGARefLink{sga3:VIIA:3.2.3}{3.2.3})}, we then have
\[
\begin{aligned}
  \Delta_U\bigl(\exp_U(x)\bigr)
  &=
  \exp_{U\otimes_CU}\bigl(\Delta_U(x)\bigr)
  \\
  &=
  \exp_{U\otimes_CU}(x\otimes1+1\otimes x)
  \\
  &=
  \exp_{U\otimes_CU}(x\otimes1)\,
  \exp_{U\otimes_CU}(1\otimes x)
  \\
  &=
  \bigl(\exp_U(x)\otimes1\bigr)
  \bigl(1\otimes\exp_U(x)\bigr)
  \\
  &=
  \exp_U(x)\otimes\exp_U(x).
\end{aligned}
\]
Thus the bijection \(\exp_U\) sends a primitive element of
\(\mathfrak rU\) to an element \(z\in1+\mathfrak rU\) satisfying
\[
  \Delta_U(z)=z\otimes z.
\]
Conversely, if \(z\) satisfies these conditions, then, on putting
\(x=\log_U(z)\), the preceding calculation shows that
\[
  \exp_{U\otimes_CU}(x\otimes1+1\otimes x)
  =
  z\otimes z
  =
  \Delta_U\bigl(\exp_U(x)\bigr)
  =
  \exp_{U\otimes_CU}\bigl(\Delta_U(x)\bigr);
\]
hence
\[
  x\otimes1+1\otimes x=\Delta_U(x).
\]
\footnote{\label{sga3:VIIB:nde:121}The editors supplied the
details of the preceding argument and added the sentence that
follows.}
Furthermore, if \(z\in1+\mathfrak rU\) satisfies
\(\Delta_U(z)=z\otimes z\), then
\[
  \varepsilon_U(z)^2=\varepsilon_U(z).
\]
Since \(\varepsilon_U(z)\) is invertible \textup{(because \(z\) is
invertible, and \(\mathfrak r\) is nilpotent)}, it follows that
\(\varepsilon_U(z)=1\).

In particular, let \(L\) be a Lie algebra flat over \(C\), take for
\(U\) the enveloping algebra \(U(L)\) of \(L\) over \(C\), and
identify \(L\) with its canonical image in \(U\). By \SGARefLink{sga3:VIIB:3.1:lemma}{Lemma~3.1},
\(L\) is the set of primitive elements of \(U\): indeed, on every
local component of \(C\), \(L\) is a free module. Consider the
\(C\)-formal group
\[
  \mathscr G(L)=\operatorname{Spf}^{*}U(L),
\]
whose covariant bialgebra is \(U=U(L)\)
\textup{(cf.\ \SGARefLink{sga3:VIIB:2.6.2:prop}{2.6.2})}. Let \(\mathfrak m\) be a maximal ideal of
\(C\), and put
\[
  \overline C=C/\mathfrak m,
  \qquad
  \overline U=U\otimes_C\overline C.
\]
Since \(\mathscr G(L)\) is infinitesimal \textup{(loc.\ cit.)}, the
unit element of \(\overline U\) is the only element
\(\overline z\in\overline U\) such
that
\[
  \varepsilon_{\overline U}(\overline z)=1,
  \qquad
  \Delta_{\overline U}(\overline z)
  =
  \overline z\otimes\overline z.
\]
It follows that every \(z\in U\) satisfying
\[
  \varepsilon_U(z)=1,
  \qquad
  \Delta_U(z)=z\otimes z
\]
necessarily belongs to \(1+\mathfrak rU\).

Finally, since
\[
  L\cap\mathfrak rU
  =
  \mathfrak rL
  =
  L\otimes_C\mathfrak r,
\]
we conclude that \(\exp_U\) defines a bijection from
\(L\otimes_C\mathfrak r\) onto the group \(\mathscr G(L)(C)\).
We summarize this as follows.

\medskip
\noindent\textbf{Proposition 3.2.}
\phantomsection
\label{sga3:VIIB:3.2:prop}
{\itshape Let \(k\) be a pseudocompact ring containing
\(\mathbb Q\), and let \(\mathbf L\) be a flat Lie algebra over
\(\mathbf O_k\).
\begin{enumerate}
\item[\textup{(i)}]
For every object \(C\) of \(\operatorname{Alf}/k\), let
\(\mathfrak r(C)\) denote its radical. Then the map
\[
  \exp_{U(\mathbf L(C))}:
  \mathbf L(C)\otimes_C\mathfrak r(C)
  \longrightarrow
  \mathscr G(\mathbf L)(C)
\]
is bijective and functorial in \(C\) and \(\mathbf L\).

\item[\textup{\SGAInlineRefTarget{sga3:VIIB:3.2:prop:item:ii}(ii)}]
The natural morphism
\[
  \mathbf L
  \longrightarrow
  \operatorname{Lie}\bigl(\mathscr G(\mathbf L)\bigr)
\]
\textup{(cf.\ \SGARefLink{sga3:VIIB:2.6.2:prop}{2.6.2})} is an isomorphism of Lie algebras over
\(\mathbf O_k\).
\end{enumerate}}
\footnote{\label{sga3:VIIB:nde:123}The editors added part
\textup{\SGARefLink{sga3:vii:paragraph:5-1:equation:ii}{(ii)}}, which follows immediately from \SGARefLink{sga3:VIIB:3.1:lemma}{3.1}.}

